
\section{Related Work}

\subsection{Efficient Attention Architectures}

\subsubsection{Sparse Attention}

Longformer (Beltagy et al., 2020), BigBird (Zaheer et al., 2020). These 
methods design sparse attention patterns but do not prune tokens. EAT 
combines standard attention with dynamic pruning.

\subsubsection{Low-Rank Attention}

Linformer (Wang et al., 2020) uses low-rank key/value projections. 
Nystr mformer (Xiong et al., 2021) uses Nystr m approximation. EAT 
preserves standard attention shape.

\subsubsection{Kernel Attention}

Performer (Choromanski et al., 2021) and Linear Transformer (Katharopoulos 
et al., 2020). EAT uses standard softmax attention.

\subsection{Token Pruning and Reduction}

\subsubsection{Vision Transformers}

DynamicViT (Rao et al., 2021) and EViT (Liang et al., 2022) prune tokens in 
vision transformers based on attention scores. EAT borrows this idea but 
applies it to text.

\subsubsection{Text Transformers}

LTP (Kim et al., 2022) uses a learned scorer to prune tokens. PuMer (Cao et 
al., 2023) uses a similar approach with vision-language models. EAT differs 
by combining pruning with LoRA recovery and multi-architecture support.

\subsubsection{Token Merging}

ToMe (Bolya et al., 2023) merges similar tokens rather than pruning. EAT 
prunes outright; the merge alternative is a future direction.

\subsection{Parameter-Efficient Fine-Tuning}

LoRA (Hu et al., 2021), AdaLoRA (Zhang et al., 2023), DoRA (Liu et al., 2024). 
EAT uses standard LoRA for recovery.

\subsection{Knowledge Transfer}

DistilBERT (Sanh et al., 2019) distills. EAT transfers weights exactly 
(zero-shot). The two are orthogonal.

\subsection{Position of Our Work}

\begin{table}[H]
\centering
\small
\begin{tabular}{lll}
\toprule
\textbf{Axis} & \textbf{Prior work} & \textbf{EAT} \\
\midrule
\textbf{Architecture} & Single family & Four families \\
\textbf{Weight transfer} & Retrain or distill & Zero-shot transfer \\
\textbf{Recovery} & Full or LoRA & LoRA only \\
\textbf{Pruning} & Train-time & Inference-time \\
\textbf{Parameter effect} & Reduces & \textbf{Increases} (LoRA) \\
\bottomrule
\end{tabular}
\end{table}

\textbf{Key distinction.} EAT is not a compressor. It replaces attention to 
support \textbf{inference-time pruning} and \textbf{multi-architecture deployment}. 
Any parameter reduction reported for EAT would be a misattribution.

\subsection{Comparison to Companion Modules}

\begin{lstlisting}

---
\end{lstlisting}